\author{OTS}
\documentclass[14pt]{../../inc/mybib}

\setincpath{../../inc/}

\usepackage{bible_style}
\usepackage{textmakro}
\graphicspath{{../../assets/images/}}
\usepackage{header}

\newcommand{\Name}{Samuel}

% ensure scrlayer-scrpage has sufficient footheight
\setlength{\footheight}{20.4pt}

\begin{document}
\begin{bibelbox}{SCHL}{Kol}{3:15-17}
    Lasst das Wort des Christus reichlich in euch wohnen in aller Weisheit; lehrt und ermahnt einander und singt mit Psalmen und Lobgesängen und geistlichen Liedern dem Herrn lieblich in eurem Herzen. Und was immer ihr tut in Wort und Werk, das tut alles im Namen des Herrn Jesus und dankt Gott, dem Vater, durch ihn.
\end{bibelbox}
\section{Begrüssung}
Ich möchte auch euch alle recht herzlich zum Gottesdienst begrüssen. Schön, dass ihr gekommen seid, um unserem \herr N zu danken, ihn zu loben und zu preisen.

Auch dich \Name{} will ich herzlich begrüssen. Schön, dass du wieder bei uns bist und mit uns Gottesdienst feierst.
\begin{itemize}
    \item[] \beten{}
    \item[] Und anschliessend singen wir zusammen das Lied
    \item[] \lied{Wunderbar. grosser Erlöser}.
\end{itemize}

\section{Informationen}
\begin{itemize}
    \item \bt[infos]{Bibel und Gebetsabend:} Do, 12.03.2026 20:00 Uhr Bibel und Gebetsabend ; zu Markus 6, 1-6 mit Elia Morris.
    \item \bt[infos]{Nächster Gottesdienst:} So, 15.03.2026 14:45 Uhr mit Fredy Peter.
    \item \bt[infos]{Kollekte:} Die Kollekte wird hinten in der grünen Box gesammelt und wird vollumfänglich für den Bau dieser Gemeinde eingesetzt.
\end{itemize}

\section{ Input }
\begin{spacing}{1.5}
    \begin{block}[Die Trübsahlzeit]
        Mit der 70. Jahrwoche in \bibleverse{Dan}(9:24-27) beginnt die Trübsahlzeit hier auf der Erde. Letzten Sonntag haben ich euch gesagt, dass ich glaube, dass die Entrückung vor dieser Zeit stattfinden wird. Das heisst also, wenn wir jetzt über die Trübsahlzeit sprechen sind wir als Gemeinde nicht mehr da. Es sind allso nur noch Menschen auf dieser Welt die Jesus Christus als ihren Retter nicht angenommen haben. Zu Beginn dieser letzten sieben Jahre, werden die Opfer im Tempel wieder eingeführt. Aber nicht Gott ist es, der diese Opfer wieder einführt, sondern der Antichrist. Zu Beginn bringt dieser eine Ruhe und Frieden. Liebi nennt diese Zeit die Stunde der Versuchung. Die Menschen auf der Erde verehren ihn. Wenn einer Auftritt und den Frieden auf Erden bringt, dann glaubt mir, jeder nicht Bibeltreue Christ wird auf diesen Menschen reinfallen. Das sieht man ja schon heute wie Menschen von der Kirche gefeiert werden die irgend etwas gutes Tun. Aber in Römer Kap. steht, dass wir aus uns nichts Gutes tun können. Es tut Gott. Ihm soll man danken. Das wir hier sitzen und uns nicht durch diese Zeit durch müssen verdanken wir Gott, der uns in seiner Gnade errettet hat.
    \end{block}
    \lied{Wir wissen nicht den Tag noch die Stunde}.
\end{spacing}
\section{Predigt}
\textbf{Nach der Predigt}\newline
Danken für die Predigt.
\section{Abendmahl}
\begin{itemize}
    \item [] Beten für das Brot
    \item [] Beten für den Wein
\end{itemize}
\section{Abschluss}
\begin{itemize}
    \item [] \lied{O Gott dir sei Ehre}.
    \item [] \beten{}
\end{itemize}
\begin{bibelbox}{SCHL}{IIKor}{13:13}
    Die Gnade des Herrn Jesus Christus und die Liebe Gottes und die Gemeinschaft des Heiligen Geistes sei mit euch allen! Amen
\end{bibelbox}
\end{document}